\documentclass[11pt,a4paper]{article}
\usepackage{amsmath,amssymb,amsthm}
\usepackage[margin=1.05in]{geometry}
\usepackage{longtable}
\usepackage{hyperref}

\theoremstyle{plain}
\newtheorem{theorem}{Theorem}
\newtheorem{conjecture}[theorem]{Conjecture}
\newtheorem{corollary}[theorem]{Corollary}
\newtheorem{proposition}[theorem]{Proposition}
\newtheorem{lemma}[theorem]{Lemma}
\theoremstyle{definition}
\newtheorem{remark}[theorem]{Remark}

\renewcommand{\SS}{\mathfrak{S}}
\newcommand{\AAA}{\mathfrak{A}}

\title{The demand side of the Elliott--Halberstam route to binary
Goldbach,\\ and the wall's exact second moment}

\author{[author]}
\date{Version 2\\
\small All measurements reproducible from the accompanying repository}

\begin{document}
\maketitle

\begin{abstract}
Huang and Li proved that binary Goldbach for large even $N$ follows
from $EH_\mu(N^{\theta'})$ for a single $\theta' > 1/2$. This paper
maps that hypothesis from both of its sides and establishes what is
exactly computable about the single scalar it reduces to.

On the \emph{demand} side we prove one unconditional theorem
(Theorem~\ref{thm:A}): the M\"obius-weighted, fixed-class correlation
sum with weight $w_k = 1$ is $\ll_A N(\log N)^{-A}$, so Huang--Li's
$E_4$ consumption of $EH_\mu$ is unnecessary and the whole demand
collapses to the single scalar $E_3$. That scalar is \emph{equivalent}
to binary Goldbach by Huang--Li's own equation~(22)
(Theorem~\ref{thm:C}), so the demand-side search for a weaker
sufficient condition is closed at the level of identities. We then
close the interior of the design space those two endpoints open
(Theorems~\ref{thm:D} and~\ref{thm:Dprime},
Propositions~\ref{prop:E} and~\ref{prop:Dpp}). We also report a defect
in equation~(18) of the source paper, and its repair.

The demand side reduces everything to
$C(N)=\sum_{n<N}\Lambda(n)\mu(N-n) = o(N)$.
Sections~\ref{sec:wall}--\ref{sec:floor} give the exact facts about
that quantity: its second moment in closed form
(Proposition~\ref{prop:V}), an exact identity for the aggregate second
moment (Lemma~\ref{lem:MP}), an exact closed form for the fluctuation
of any cell mean (Lemma~\ref{lem:cellmom}), and the observation that
this error bar falls like $(\log N)^{-1/2}$ and not like $n_c^{-1/2}$
(Proposition~\ref{prop:coh}), which changes the width of every interval
in this subject built from a count. Two controls
(Lemmas~\ref{lem:coin} and~\ref{lem:placebo}) delimit what can be
measured at all: in particular the excess of $\mathrm{Var}\,C$ over its
random-sign value cannot be estimated by any centred statistic, because
such a statistic returns the same value on a coin.

We also quantify the obstruction to the natural conditional route:
Chowla's conjecture does \emph{not} give $\mathrm{Var}\,C \sim V$,
because the coefficient amplifying $S(h)$ grows like $N/\log N$
(Proposition~\ref{prop:W}), and since that coefficient is nonnegative
no bound on $|S(h)|$ suffices at all.

The deliverable is the map and the exact facts, not a theorem toward
Goldbach. Net progress toward Goldbach: zero. We state throughout what
is and is not claimed, and \S\ref{sec:notclaimed} states what has been
deliberately withheld from this version.
\end{abstract}

\tableofcontents

\section{Introduction}

\subsection{The conditional frame}

Let $N$ be a large even integer,
$\tilde r(N) = \sum_{n<N}\Lambda(n)\Lambda(N-n)$, and $\SS(N)$ the
singular series. Write $EH_\mu(Q)$ for the assertion that for every
$A>0$
\[
  \sum_{q\le Q}\max_{y<N}\max_{(a,q)=1}
  \Bigl|\sum_{\substack{n\le y\\ n\equiv a\,(q)}}\Lambda(n)\mu(N-n)
   -\frac{1}{\varphi(q)}\sum_{n\le y}\Lambda(n)\mu(N-n)\Bigr|
  \ll_A \frac{N}{(\log N)^A}.
\]
Huang and Li~\cite{HL}, continuing Pan~\cite{Pan}, proved that
$EH(N^{\theta}(\log N)^{2A+8})$ together with $EH_\mu(N^{1-\theta})$
implies $\tilde r(N)\ge \SS(N)(1-A(N))N + O(N(\log N)^{-A})$, and
deduced (their Corollary~1) that in view of Bombieri--Vinogradov,
binary Goldbach for all sufficiently large even integers follows from
$EH_\mu(N^{\theta'})$ alone for any single $\theta' > 1/2$.

This is the sharpest conditional frame we are aware of: one sentence
stands between classical technology and Goldbach. This paper asks what
that sentence actually costs.

\subsection{The organizing principle: two couplings}

The hypothesis contains a product of two M\"obius factors, and there
are exactly two ways their arguments can be coupled. This turns out to
organize everything.

\begin{itemize}
\item \textbf{Divisibility-coupled} (the \emph{demand} side). Inside
  Huang--Li's proof, $EH_\mu$ is consumed only in the fixed residue
  class $n \equiv N \pmod k$, i.e.\ $k \mid N-n$: one M\"obius
  argument divides the other. Section~\ref{sec:demand}.
\item \textbf{Difference-coupled} (the \emph{supply} side). Any attempt
  to \emph{prove} $EH_\mu$ by dispersion arrives at the dilate field
  $D(k) = \sum_m \mu(m)\mu(N-mk)$, where the two arguments are related
  by a difference, not a divisibility. Section~\ref{sec:supply}.
\end{itemize}

\noindent
The two behave completely differently, and each is closed for its own
reason:

\begin{center}
\begin{tabular}{p{0.20\textwidth}p{0.34\textwidth}p{0.34\textwidth}}
\hline
 & \textbf{Divisibility-coupled} & \textbf{Difference-coupled}\\
\hline
Divisor switch & exact, and free & exact, but useless\\
M\"obius lands on & the \emph{short} variable, $m\le N^{1/2-\delta}$ &
 the \emph{long} variable, $m \asymp N/K \ge N^{2/3}$\\
Consequence & Bombieri--Vinogradov closes it: Theorem~\ref{thm:A} &
 no known machine applies\\
But & the weighted version is \emph{equivalent} to Goldbach
 (Theorem~\ref{thm:C}) & no named coupling surface found, internal or
 external (\S\ref{sec:closures})\\
\hline
\end{tabular}
\end{center}

\subsection{What this paper claims}\label{sec:claims}

\begin{itemize}
\item \textbf{Claimed}: Theorem~\ref{thm:A}, Corollary~\ref{cor:B},
  Theorem~\ref{thm:C}, Theorems~\ref{thm:D} and~\ref{thm:Dprime},
  Propositions~\ref{prop:E}, \ref{prop:Dpp}, \ref{prop:V},
  \ref{prop:W}, \ref{prop:scaleinv} and~\ref{prop:coh},
  Lemmas~\ref{lem:MP}, \ref{lem:cellmom}, \ref{lem:coin}
  and~\ref{lem:placebo}; the defect in \cite{HL}'s equation~(18) and
  its repair; the measurements reported here, which are reproducible.
\item \textbf{Conjectured}: Conjecture~\ref{conj:L}.
\item \textbf{Not claimed}: any theorem toward Goldbach.
  Theorem~\ref{thm:A} is unconditional but Goldbach-neutral: it
  removes exactly the half of the demand that carries no Goldbach
  content. See also \S\ref{sec:notclaimed}.
\end{itemize}

\section{The demand side}\label{sec:demand}

\subsection{Where $EH_\mu$ is actually consumed}

For $(k,N)=1$ and $1\le t<N$ put
\[
  E_\mu(t;k) = \sum_{\substack{n\le t\\ n\equiv N\,(k)}}
    \Lambda(n)\mu(N-n) - \frac{1}{\varphi(k)}\sum_{n\le t}
    \Lambda(n)\mu(N-n).
\]
With $K = (N-1)/\alpha$, Huang--Li's two consumption sites are
\[
  E_3(\alpha)=\!\!\sum_{\substack{k<K\\(k,N)=1}}\!\!\mu(k)\log k\,
  E_\mu(N;k),
  \qquad
  E_4(\alpha)=\!\!\sum_{\substack{k<K\\(k,N)=1}}\!\!\mu(k)\,
  \bigl[\text{$E_\mu$ with the extra weight }\log(N-n)\bigr].
\]
They differ in exactly one respect: the weight $w_k$ is $\log k$ in
$E_3$ and $1$ in $E_4$ (the $\log(N-n)$ inside $E_4$ is
$k$-independent and comes off by partial summation). In both cases
Huang--Li discard the signs $\mu(k)$ on the first line by the triangle
inequality, and only then appeal to $EH_\mu$. Keeping the signs is
what the following results exploit.

\subsection{The theorem}

\begin{theorem}[$w_k = 1$]\label{thm:A}
Fix $\theta' \in (1/2,1)$ and put $K = N^{\theta'}$. Then for every
$A>0$,
\[
  \sup_{1\le t<N}\Bigl|
   \sum_{\substack{k<K\\(k,N)=1}}\mu(k)\,E_\mu(t;k)\Bigr|
  \;\ll_{A,\theta'}\; \frac{N}{(\log N)^A}.
\]
Every ingredient except Bombieri--Vinogradov contributes only
$O(Ne^{-c\sqrt{\log N}})$.
\end{theorem}

The mechanism is short, but it has three steps and not one. The class
$n \equiv N \pmod k$ is $k \mid N-n$, so with $u = N-n$ the $k$-sum is
the incomplete divisor sum $\sigma_K(u) = \sum_{k\mid u,\,k<K}\mu(k)$.
\emph{First} one completes it: for squarefree $u$,
$\sum_{k\mid u,\,(k,N)=1}\mu(k) = \mathbf 1_{\mathrm{rad}(u)\mid N}$,
which forces $u\mid \mathrm{rad}(N)$ and so contributes $N^{o(1)}$
terms. \emph{Then} the complementary sum, over $k\ge K$, is the one
that carries the work --- and it is only there that writing $u = mk$
gives
\[
  m \;=\; u/k \;<\; N/K \;=\; N^{1-\theta'} \;\le\; N^{1/2-\delta}.
\]
\emph{Finally}, for squarefree $u=mk$ one has $(m,k)=1$ and
$\mu(u)\mu(k) = \mu(m)\mu^2(k)$: the M\"obius on the long variable $k$
squares itself away and the surviving M\"obius sits on the short
variable $m$. That is the classical ``M\"obius on the short variable
plus Bombieri--Vinogradov'' configuration.

The completion step is not decorative. Over the range $k<K$ as it
stands, the cofactor $u/k$ runs up to $u$ itself, and the M\"obius
would sit on the \emph{long} variable --- the configuration the table
of \S1.2 records as having no known machine. Full proofs are in the
companion note~\cite{ThmA}.

The bound is $N(\log N)^{-A}$ for each fixed $A$ and no better. It is
not $Ne^{-c\sqrt{\log N}}$: Bombieri--Vinogradov's internal
Siegel--Walfisz range is $q\le(\log N)^B$ with $B$ fixed, and that is
what caps the saving.

\begin{corollary}\label{cor:B}
$E_4(\alpha) \ll_A N(\log N)^{-A}$ unconditionally, via the exact
identity $E_4(\alpha) = \int_1^{N-1} T_1(t)\,dt/(N-t)$, where
$T_1(t)$ is the sum of Theorem~\ref{thm:A}. Hence Huang--Li's Lemma~4
is not needed and the entire $EH_\mu$ demand collapses to the single
scalar $E_3(\alpha)$.
\end{corollary}

\begin{proof}[Derivation]
$\int_n^{N-1}dt/(N-t) = \log(N-n)$, so for any coefficients $a_n$,
$\sum_n a_n\log(N-n) = \int_1^{N-1}\bigl(\sum_{n\le t}a_n\bigr)
dt/(N-t)$. Applying this with $a_n$ the summands of $T_1$ and bounding
$|E_4|\le \sup_t|T_1(t)|\cdot\log(N-1)$ gives the corollary, $A$ being
arbitrary.
\end{proof}

\begin{theorem}[$w_k = \log k$: the permanent closure]\label{thm:C}
Unconditionally,
\[
  E_3(\alpha) = \sum_{n<N}\Lambda(n)\Lambda(N-n)
   - \SS(N)\Bigl(N - \sum_{n<N}\Lambda(n)\mu(N-n)\Bigr)
   + O_A\!\left(N(\log N)^{-A}\right),
\]
which is Huang--Li's own equation~(22). Hence
$E_3(\alpha)\ll_A N(\log N)^{-A}$ --- the weakest form of the
hypothesis their argument needs --- is \emph{equivalent} to
$\tilde r(N)\sim\SS(N)N$.
\end{theorem}

Theorem~\ref{thm:C} is not new: \cite{HL} state the equivalence
themselves, immediately after their equation~(22). It is recorded here
because it is the other endpoint of the design space that
Theorem~\ref{thm:A} opens, and because the next subsection closes the
interior between them.

\begin{remark}
The two theorems are the same computation with two weights. For
$w_k = \log k$ the complete divisor sum is $-\Lambda$ rather than an
indicator, because $\mu * \log = \Lambda$ --- and that identity is
where Huang--Li \emph{start} (their~(10)). The $\log k$ weight is the
carrier of the Goldbach content, so every divisor switch must hand it
back. This is closure at the level of identities, not of estimates: no
choice of $\theta'$, truncation, or smoothing can evade it.
\end{remark}

\subsection{The interior of the weight space is empty}

Theorems~\ref{thm:A} and~\ref{thm:C} are the two endpoints of a design
space: the weight $w_k$ attached to $\mu(k)$ in the demand functional.
It is natural to look for a weight between them --- complete divisor
sum still cheap, main-term coefficient still of size $1$ --- that would
extract $C(N)=\sum_{n<N}\Lambda(n)\mu(N-n)$ from
Bombieri--Vinogradov alone. There is none.

Write $b := \mu*w$, so $w_k=\sum_{d\mid k}b_d$; every weight has this
form. For squarefree $u$ one has $\sum_{k\mid u}\mu(k)w_k=\mu(u)b_u$,
so the switch identity reads
$B_w\,C(N)=\sum_{u<N}\Lambda(N-u)\mu^2(u)b_u-\mathcal R_w$ with
$B_w=\sum_{k<K,(k,N)=1}\mu(k)w_k/\varphi(k)$. Two exact facts then
oppose each other:

\begin{itemize}
\item \emph{Extraction}:
  $B_w=\sum_{d<K}b_d\,\frac{\mu(d)}{\varphi(d)}\rho_{dN}(K/d)$ with
  $\rho \ll e^{-c\sqrt{\log x}}$, so $B_w$ is controlled by the part of
  $b$ sitting \emph{at} the truncation point $K=N^{\theta'}$.
\item \emph{BV-accessibility}: expanding $w$ inside the residual gives
  moduli $md$ with $m<N^{1-\theta'}$, so $b$ may only be supported on
  $d\le N^{\theta'-1/2-\delta}$.
\end{itemize}

\begin{theorem}[no weight extracts $C(N)$]\label{thm:D}
Fix $\delta>0$. If $b=\mu*w$ is supported in
$[1,N^{\theta'-1/2-\delta}]$ then
$\|b\|_1/|B_w| \gg \exp\bigl(c\sqrt{(1/2+\delta)\log N}\bigr)$, and the
switch identity yields at best
\[
  |C(N)|\;\ll\;\exp\bigl(c\sqrt{(\tfrac12+\delta)\log N}\bigr)\,
  N(\log N)^{-A},
\]
which is no saving of any power of $\log N$. Letting $\delta\to0$ the
loss exponent tends to $\sqrt{\tfrac12\log N}$ and never below it.
\end{theorem}

\begin{theorem}[the no-go survives $EH$]\label{thm:Dprime}
If $\Lambda$ has level of distribution $\theta_E<1$, the same argument
gives $\|b\|_1/|B_w| \gg \exp(c\sqrt{(1-\theta_E)\log N})$. Closing the
gap requires $\theta_E=1$: equidistribution of $\Lambda$ to moduli of
size $N$ itself, where each progression holds $O(1)$ terms.
\end{theorem}

\begin{proposition}[the circle method has zero margin]\label{prop:E}
Writing $C(N)=\int_0^1 S_\Lambda S_\mu e(-N\alpha)d\alpha$, both
standard estimates lie at or above the trivial bound $\psi(N)\sim N$.
Cauchy--Schwarz gives
$\|S_\Lambda\|_2\|S_\mu\|_2 \sim (6/\pi^2)^{1/2}N(\log N)^{1/2}$,
exceeding it by a growing factor. For the pointwise route, Parseval
gives $\sup_\alpha|S_\mu| \ge \|S_\mu\|_2 = (6N/\pi^2)^{1/2}$, so no
improvement in M\"obius exponential-sum technology can push the first
factor below the scale $N^{1/2}$; and $\|S_\Lambda\|_1 \gg N^{1/2}$ by
Vaughan~\cite{Vau88}, so the product is $\gg N$.
The measured margin $N/(\sup|S_\mu|\,\|S_\Lambda\|_1)$ is
$0.1679,\,0.1749,\,0.1576,\,0.1521$ at
$N=2^{14},2^{16},2^{18},2^{20}$ --- below $1$ throughout, and falling
over the last three.
\end{proposition}

Davenport's uniform bound $S_\mu\ll_A N(\log N)^{-A}$ is useless here:
it saves against $N$, while the pairing needs the scale $N^{1/2}$,
which Parseval forbids. This is the parity obstruction in
circle-method language; compare \cite{MV17} for the same barrier stated
in sieve terms.

Theorem~\ref{thm:D} assumes $b$ is supported low enough for BV, which
excludes the most natural family: $w_k=f(\log k)$ with $f$ polynomial.
There $b=\mu*\log^D=\Lambda_D$, the generalised von Mangoldt function,
so the complete part splits by the number of prime factors of $u$ into
a Goldbach-type piece ($r=1$) and Chen-type pieces ($r\ge2$).

\begin{proposition}[polynomial weights]\label{prop:Dpp}
$D=0$ gives $B_w\ll e^{-c\sqrt{\log K}}$. For a monomial $x^D$,
$D\ge1$, every term of the complete part is nonnegative
($\Lambda\ge0$, $\Lambda_D\ge0$), so it is $\asymp N(\log N)^{D-1}$
with a fixed sign. Cancelling across monomials requires tuning against
the asymptotics of the $r=1$ piece, which at $D=1$ is the binary
Goldbach sum --- circular. No polynomial weight makes the complete
part a classically known object.
\end{proposition}

Measured at $N=10^6,4\cdot10^6,1.6\cdot10^7$: the two pieces of
$\mathrm{CP}_2$ are the same order, with ratio
$0.7712,\,0.7903,\,0.8103$, and
$\mathrm{CP}_2/(N\log N) = 2.8864,\,2.9493,\,2.9966$. The closure rests
on nonnegativity, not on one piece dominating; the drift of the ratio
is real over the range computed but reaching $1$ at its observed rate
would take $N$ near $10^{13}$, so nothing is claimed about its limit.
The one analytically canonical tuning, $f(x)=x^2-2\gamma x$, which
makes $b$ mean-zero by killing the pole of $\zeta W$ at $s=1$, moves
the complete part by $-5.1\%$, $-4.5\%$, $-4.1\%$ at the three $N$:
mean-zero is not smallness, because it removes the average of $b$ and
not its correlation with $\Lambda(N-\cdot)$.

The loss exponent $\sqrt{(1/2)\log N}$ is the $\sqrt N$ barrier itself:
it is the gap between where a weight must live to see $C(N)$ and where
Bombieri--Vinogradov permits it to live. This is a no-go for one
precisely specified method --- divisor switching with BV as the only
input --- over that method's entire weight space; it is not an
obstruction to other methods. Its purpose is to close the design space
that Theorems~\ref{thm:A} and~\ref{thm:C} opened.

\subsection{A defect in the published paper}

Equation~(18) of \cite{HL} replaces the $n$-dependent constraint
$k < (N-n)/\alpha$, present three displays earlier in the derivation of
$S_2(\alpha)$, by the $n$-free $k < (N-1)/\alpha$. The terms thereby
included are exactly those with $d=(N-n)/k \le \alpha$, and they sum to
\[
  \Delta = -\sum_{2\le m\le\alpha}\mu(m)\log m
    \sum_{\substack{k<(N-1)/\alpha\\ (k,m)=1}}\mu^2(k)\Lambda(N-mk),
\]
whose trivial bound $\ll N(\log N)^2$ exceeds the target. It closes by
the same machinery --- M\"obius again on the short variable --- under
hypotheses already assumed: unconditionally by Bombieri--Vinogradov in
the Corollary-1 regime, and under the assumed $EH$ in the Theorem-1
regime. Theorem~1 and Corollary~1 of \cite{HL} stand as stated.

\section{The supply side}\label{sec:supply}

\subsection{The consumable}

The chain's supply-side consumable, distilled, is
\[
  \textbf{E1 (weak form).}\quad
  D(k) = \!\!\sum_{\sqrt N < m \le N/k}\!\! \mu(m)\mu(N-mk),
  \qquad
  \sum_{k\sim K}|D(k)|^2 \ll (\log N)^{-2A-2}\sum_{k\sim K} M_k^2,
\]
for $K \le N^{1/3}$ dyadic, where $M_k$ is the length of the $m$-range.
A fixed log-power saving is the currency; $o(1)$ and $\log\log$ savings
are not consumable.

The normalisation is $M_k^2$, i.e.\ the \emph{trivial} bound, not
$M_k$, the square-root scale: the wall is
$T_{II} = \sum_{k\sim K}b_kD(k)$ with $|b_k|\ll\log N$, needed at
$N(\log N)^{-A}$, and Cauchy--Schwarz in $k$ gives
$\sum_k|D(k)|^2 \ll N^2/(K(\log N)^{2A+2})$ with
$\sum_{k\sim K}M_k^2 \asymp N^2/K$. The distinction is load-bearing:
at the square-root normalisation the demand would be
$\sum|D|^2/\sum M_k \ll 8.8\cdot10^{-6}$ at $N=10^8$, $A=1$, which our
own measurements refute directly --- the measured ratio is
$0.34,\,0.39,\,0.32$ over three dyadic bands, since
$\prod_p(1-2/p^2) = 0.32263$ is the density of $m$ for which both
$\mu(m)$ and $\mu(N-mk)$ are nonzero. Normalised by that surviving
support the ratio is $1.04,\,1.17,\,0.96$: the field exhibits
\emph{exact} square-root cancellation on the terms that survive. At the
correct normalisation the target is cleared with margin
$(N/K)(\log N)^{-2A-2}$, which is asymptotic and is below $1$
everywhere accessible, so no computation is offered as evidence for it.

\begin{remark}[the mandatory comparison]\label{rem:scale}
A target must be checked against one's own measurement of the same
quantity before it is used to adjudicate anything. Writing a target at
the square-root scale when the chain consumes it at the trivial scale
inflates every demand by a power of $N$, and no verdict passed on a
budget basis survives it. We record this as a standing rule because
the same slip occurred twice in the course of this work, and both times
it was caught only by that comparison.
\end{remark}

\subsection{Conjecture L}

\begin{conjecture}[the factorization law]\label{conj:L}
Each M\"obius family probed in this program --- prime-indexed dilate
pairs $C_{k,k'} = \sum_{p\sim P}\mu(N-pk)\mu(N-pk')$, integer-indexed
dilates and their pairs, and M\"obius sums over thin progressions ---
factorizes as
\[
  \text{field} \;=\; \mathbf M \times \mathbf G,
\]
where $\mathbf M$ is the deterministic local mask, computed exactly by
finite modular enumeration from the $v_q$-data of $(N,k,k')$, and
$\mathbf G$ is, on the surviving support, a fluctuation with no mean
field and no class structure.
\end{conjecture}

The evidence is of two kinds and we separate them. The E1 arm is stated
above and has been reproduced independently. The remaining stamps ---
pair statistics, the exact $(v_2,v_3)$ cells, the blind mask
prediction, and the Wishart pair spectrum --- are recorded in the
repository and have \emph{not} been independently re-verified; they are
cited here as the reason the conjecture is stated, not as measurements
this paper vouches for.

Conjecture~\ref{conj:L} is \emph{stronger} than the chain needs: E1
requires only square-root cancellation on average. What is missing is
therefore not knowledge of the structure --- the mask is an algorithm
and $\mathbf G$ is featureless in every measurement --- but a proof
technique for the amplitude of a featureless object.

\section{The wall, exactly}\label{sec:wall}

The chain of Section~\ref{sec:demand} reduces the whole problem to one
scalar,
\[
  C(N) \;=\; \sum_{n<N}\Lambda(n)\mu(N-n) \;=\; o(N).
\]
This section gives that scalar its exact second moment, an exact
identity for its aggregate second moment, and the two controls that
delimit what can be measured about it at all.

\subsection{The exact scale}

\begin{proposition}[The exact scale]\label{prop:V}
Let
\[
  V(N) \;=\; \sum_{v<N}\mu^2(v)\,\Lambda(N-v)^2,
  \qquad
  W(N) \;=\; \sum_{w<N}\Lambda(w)^2,
\]
and let
\[
  \AAA(N) \;=\; \prod_{q\,\nmid\,N}\Bigl(1-\frac{1}{q(q-1)}\Bigr).
\]
Then $V(N) = W(N)\,\AAA(N)\,(1+o(1))$, and consequently
$V(N) \sim \AAA(N)\,N\log N$.
\end{proposition}

\begin{proof}[Derivation]
$\Lambda(w)^2$ is supported on prime powers with weight $(\log p)^2$,
so with $w=N-v$,
\[
  V(N) \;=\; \sum_{p^k<N}(\log p)^2\mu^2(N-p^k)
        \;=\; \sum_{p<N}(\log p)^2\mu^2(N-p) + O(\sqrt N\log^2 N),
\]
a $(\log p)^2$-weighted count of \emph{squarefree shifted primes}. Its
local density at a prime $q$ is $1$ when $q\mid N$ --- there
$q^2\mid N-p$ forces $q\mid p$, hence $p=q$, a single term --- and
$1-1/\varphi(q^2) = 1-1/(q(q-1))$ when $q\nmid N$, since the bad class
$p\equiv N \pmod{q^2}$ is then a unit class. The statement is Mirsky's
theorem on squarefree values of shifted primes \cite{Mir49} together
with partial summation; it is recalled here, not claimed.
\end{proof}

\paragraph{The local factor is $\AAA$, not $\SS$.}
The singular series $\SS(N)$ is Hardy--Littlewood's and is correct for
the Goldbach \emph{count}; the right object for the second moment of
$C(N)$ is $\AAA(N)$. Both are products over the primes dividing $N$,
which is why the substitution is easy to make and hard to see.
Rescaling each candidate to the measured mean, so that only its shape
in $N$ is judged, the residual standard deviation over even $N$ in
$[10^5,\,1.6\cdot10^7]$ is $0.000323$ for $\AAA$ against $0.245235$ for
$\SS$ --- a factor of $759$. Dividing through by $W(N)$ removes the
analytic factor exactly, so no asymptotic for $\sum(\log p)^2$ is
assumed: the ratio $V(N)/W(N)$ against $\AAA(N)$ has mean $1.000000$
with standard deviation $0.000166$ in the top octave, and agrees to six
decimals in each of the six radical cells from $2\mid N$ through
$2\cdot3\cdot5\cdot7\cdot11\cdot13\mid N$. At $N=4\cdot10^6$ the ratio
$W/V$ is $1.27080$ against the predicted $1/\AAA(N) = 1.27020$.

The lower cutoff $10^5$ is not cosmetic. Below it the residual is
dominated by small-$N$ noise and the figure for $\AAA$ degrades to
$0.000582$; the figure for $\SS$ is unaffected, because it is dominated
by $\SS$'s wrong shape rather than by noise. Any comparison of the two
must state its range.

\subsection{The aggregate second moment, exactly}

\begin{lemma}[the second-moment identity]\label{lem:MP}
Fix $X$ and let
\[
  \widehat C(N) \;=\; \sum_{\substack{n+v=N\\ n\le X,\ v\le X}}
    \Lambda(n)\,\mu(v)
\]
be the truncated convolution, whose support runs to $2X$. Then
\[
  \sum_{N\le 2X} \widehat C(N)^2 \;=\; \sum_{|h|<X} M(h)\,P(h),
  \qquad
  M(h) = \sum_{v,\,v+h\le X} \mu(v)\mu(v+h),
  \quad
  P(h) = \sum_{w,\,w+h\le X} \Lambda(w)\Lambda(w+h).
\]
\end{lemma}

\begin{proof}
Expanding the square as a double sum over $(n,v)$ and $(n',v')$ with
$n+v=n'+v'$ and summing over \emph{all} $N$ leaves the single
constraint $n-n' = v'-v =: h$, with $(n,n')$ and $(v,v')$ otherwise
ranging independently over $[1,X]^2$. Collecting the $\Lambda$-pairs
and the $\mu$-pairs separately gives the two factors.
\end{proof}

\begin{remark}
The truncation on the left is necessary and it is easy to lose. If one
writes instead $\sum_{N\le X}C(N)^2$ with $C$ the untruncated
convolution, the three surviving indices $(n,v,h)$ are constrained by
$n+v\le X$ --- a \emph{simplex} --- while $M(h)P(h)$ ranges them over
the \emph{box} $[1,X]^2$ and so also counts every pair with
$X< n+v\le 2X$. The two sides then differ by a factor near $1.57$ at
$X$ from $800$ to $3200$, and the ratio does not tend to $1$. The form
above is exact to machine precision at every $X$ tested, to
$X=1.6\cdot10^7$.
\end{remark}

The identity is exact and unconditional, and it makes the aggregate
size of the wall a statement about two shifted correlations, one of
which ($P$) is the Hardy--Littlewood prime-pair sum and the other of
which ($M$) is the binary Chowla sum.

\subsection{What the excess is, and why Chowla does not control it}

Write $\rho(N) = \mathrm{Var}\,C(N)/V(N)$, which is $1$ if the $\mu(v)$
on the surviving support are replaced by independent signs. Pointwise,
$C(N)^2 = V(N) + \mathrm{OffDiag}(N)$ with
\[
  \mathrm{OffDiag}(N)
  = \sum_{h\ne0}\ \sum_{\substack{p,\,p+h\ \mathrm{prime}}}
    (\log p)(\log(p+h))\,\mu(N-p)\mu(N-p-h),
\]
which is algebra. Passing from there to $\sum_{h\ne0}c(h)S(h)$ with
$c(h)=\sum_{p'-p=h}(\log p)(\log p')$ and
$S(h)=\langle\mu(u)\mu(u-h)\rangle$ is \emph{not} algebra: it replaces
the values $\mu(N-p)\mu(N-p-h)$, sampled on the sparse set of $p$ with
$p$ and $p+h$ both prime, by an average over all $u$. We do not assume
that decoupling. What we do record is that even granting it, the
conditional route does not close.

\begin{proposition}[the amplification]\label{prop:W}
Let $\Gamma(N) = \bigl(\sum_{h\ne0}c(h)\bigr)/V(N)$. Then
$\Gamma(N) \sim N/(\AAA(N)\log N)$, and consequently a hypothesis
$|S(h)|\le\varepsilon$ yields nothing better than
$|\rho-1| \le \varepsilon\,\Gamma(N)$. Since $c(h)\ge0$, no bound on
$|S(h)|$ alone can give $\rho\to1$.
\end{proposition}

Measured, $\Gamma = 1.5489\cdot10^{3},\ 1.8517\cdot10^{4},\
3.5798\cdot10^{5}$ at $N=10^4,\,1.6\cdot10^5,\,4\cdot10^6$, with
$\Gamma\log N/N = 1.4266,\,1.3868,\,1.3605$ against the predicted
$1/\AAA \to 1.270$.

Two consequences. First, Chowla's conjecture asserts $S(h)=o(1)$ for
fixed $h$; the strength the display needs is $S(h) = o(\log N/N)$, so
the gap is a factor $N/\log N$. Second, and independently, since
$c(h)\ge0$ the only route to $\rho\to1$ is \emph{signed cancellation
across $h$}, which is a different object from Chowla's smallness and
from the averaged absolute bound of \cite{MRT15}. Naming what supplies
that cancellation is an open question and is stated as such.

\subsection{Two controls, and what they invalidate}

\begin{lemma}[the coin control]\label{lem:coin}
Let $\varepsilon(v)=\pm1$ be arbitrary signs on $\{v : \mu(v)\neq0\}$
and zero elsewhere. Then $\varepsilon^2 = \mu^2$ pointwise, so
$V(N)$ is unchanged, and the field
$C_\varepsilon(N)=\sum_v\varepsilon(v)\Lambda(N-v)$ has the same exact
second moment as $C(N)$ for every $N$. Consequently any estimator whose
output is reproduced when $\mu$ is replaced by $\varepsilon$ is not
measuring $\mu$.
\end{lemma}

\begin{lemma}[the placebo key]\label{lem:placebo}
Let the cells be indexed by a labelling $\ell(N)$, and let $\pi$ be a
permutation of the even integers in the band. Replacing $\ell$ by
$\ell\circ\pi$ preserves every cell size and leaves the field
$Z(N)=C(N)/\sqrt{V(N)}$ byte-identical, while destroying the
correspondence between cells and arithmetic. Any cell-indexed statistic
that survives this replacement is a property of the cell sizes and not
of that correspondence.
\end{lemma}

These are trivial and they are load-bearing. Lemma~\ref{lem:coin}
invalidates the natural measurement of $\rho$: the centred estimator
returns the same value on $\varepsilon$ as on $\mu$, so a quoted level
for $\rho$ calibrates nothing. \textbf{Any claim about
$\mathrm{Var}\,C$ relative to $V$ must first exhibit an estimator that
distinguishes $\mu$ from a coin, and show that it does.} With one
realisation of $\mu$ that may not be possible at all, and saying so
would itself be a result. Lemma~\ref{lem:placebo} is what confirms, in
Section~\ref{sec:floor}, that the cell floor is a property of the
arithmetic and not of the partition.

\section{The cell floor}\label{sec:floor}

Throughout, cells are indexed by \emph{depth} $d$, the number of
$3,5,7,11,13$ dividing $N$, and $Z(N)=C(N)/\sqrt{V(N)}$.

\subsection{The floor, in closed form}

\begin{lemma}[exact cell moments]\label{lem:cellmom}
Let a band of even $N$ be partitioned into cells, let $c$ be a cell of
size $n_c$ inside a band of size $n$, let $m_c$ and $\overline m$ be
the means of $Z$ over $c$ and over the band, and set
\[
  u_c(v) \;=\; \sum_{N\in c}\frac{\Lambda(N-v)}{\sqrt{V(N)}},
  \qquad
  Q_{cd} \;=\; \sum_v \mu^2(v)\,u_c(v)\,u_d(v).
\]
Then, when the $\mu(v)$ on the surviving support are replaced by
independent signs,
\[
  \mathrm{Var}\bigl(m_c-\overline m\bigr)
  \;=\; \frac{Q_{cc}}{n_c^{2}}
  \;-\; \frac{2\,Q_{ca}}{n_c\,n}
  \;+\; \frac{Q_{aa}}{n^{2}},
\]
where $a$ denotes the whole band. Every term is computable exactly by
one convolution; no simulation is involved.
\end{lemma}

\begin{proof}
Under independent signs $E[\varepsilon(v)\varepsilon(v')]
=\mu^2(v)\delta_{vv'}$, so
$\mathrm{Cov}(n_c m_c, n_d m_d) = Q_{cd}$ exactly. Expanding
$\mathrm{Var}(m_c-\overline m)$ bilinearly gives the three terms.
\end{proof}

The lemma has been checked against simulation even though it does not
need it: over $60$ independent sign draws in the band
$(10^5,2\cdot10^5]$, the Monte-Carlo standard deviation of
$m_c-\overline m$ agrees with the closed form at every depth, the
ratios running $0.88$ to $0.98$ against a Monte-Carlo precision of
$\pm0.09$.

\textbf{All three terms are needed.} The substitution
$u_c(v)\approx n_c/\sqrt V$, which gives the right scale for
$Q_{cc}/n_c^2$, is not specific to $c$; applied to all three terms it
returns the same value $S$ for each and hence $S-2S+S=0$. So it cannot
be used to estimate the variance of the difference. Exactly, the
variance is smaller than $Q_{cc}/n_c^2$ by a factor that is a
structural constant rather than noise: measured at the top octave of
$1.6\cdot10^7$, $\mathrm{Var}/(Q_{cc}/n_c^2)$ is
$0.113,\,0.016,\,0.289,\,0.551$ at depths $0,1,3,5$, reproducing to
three digits the same ratios at a quarter of that $N$.

\subsection{The floor's uncertainty does not fall like a count}

\begin{proposition}[coherent sums]\label{prop:coh}
The error bar of Lemma~\ref{lem:cellmom} does \emph{not} decay like
$n_c^{-1/2}$. For fixed $v$, about $n_c/\log N$ of the terms of
$u_c(v)$ are nonzero, each of size $\log N/\sqrt V$, so
$u_c(v)\approx n_c/\sqrt V$ and
\[
  \frac{Q_{cc}}{n_c^2}
  \;\approx\; \frac{\sum_v\mu^2(v)}{V}
  \;\asymp\; \frac{1}{\log N},
\]
independently of $n_c$; and by the previous subsection the variance
itself is a fixed fraction of this. So the standard error falls like
$(\log N)^{-1/2}$.
\end{proposition}

The heuristic is quantitatively loose --- measured
$Q_{cc}/n_c^2$ runs $0.124$ to $0.307$ by depth at
$N=1.6\cdot10^7$ against the heuristic's $0.049$ --- so it is offered
for the \emph{form} and not the constant. The form is confirmed
directly. Fitting $\mathrm{se}\propto N^{-b}$ to the exact floor across
eight octaves gives $b = 0.0395,\,0.0397,\,0.0394$ at depths $2,1,0$,
against the apparent exponent $1/(2\langle\log N\rangle) = 0.0358$ that
a $(\log N)^{-1/2}$ law predicts over that range.

\begin{remark}[a count is not an error bar]
The consequence is not confined to this paper. Over the factor $140$ in
$N$ measured here, $n_c^{-1/2}$ says the error bar shrinks by $11.8$;
it shrinks by $1.21$. \textbf{An interval built from a count is
therefore about ten times too narrow at the top of this range relative
to the bottom}, and in absolute terms the discrepancy is larger still:
at the top octave the exact floor exceeds
$\mathrm{sd}(Z)/\sqrt{n_c}$ by factors of $5.8$ to $160$, growing with
the cell, exactly as $Q_{cc}/n_c^2 \asymp 1/\log N$ predicts. The cause
is that $u_c(v)$ is a coherent sum, not a self-averaging one, and the
same is true of any cell mean of a field whose summands share a common
arithmetic input.
\end{remark}

\subsection{The mask exists}

Against the exactly computed floor of Lemma~\ref{lem:cellmom}, the cell
means are not zero. Over every octave from $6.25\cdot10^4$ to
$1.6\cdot10^7$, $\max_c|z_c|$ runs $9.1$ to $13.0$ and clears
Bonferroni in each. The signal is carried by the deep cells: at the top
octave the depths $3,4,5$ sit at $z = -1.6,\,-4.5,\,-9.1$ while depths
$0,1,2$ sit at $+0.0,\,+0.7,\,-0.1$. Under the permutation of
Lemma~\ref{lem:placebo} the floor collapses to the independent-sign
value, so what is being detected is the correspondence between cells
and divisibility and not the cell sizes.

\begin{proposition}[the size mechanism is scale-invariant]
\label{prop:scaleinv}
$D_c := E_{\mathrm{same},c}[\SS_2] - E_{\mathrm{all}}[\SS_2]$ depends
only on the distribution of the shift $h$ in the residue classes mod
$3,5,7,11,13$ that the cell fixes. That distribution is the same at
every scale, so $D_c$ is scale-invariant and predicts an exponent of
zero at every depth.
\end{proposition}

\textbf{How the mask decays is open, and this version does not answer
it.} The amplitudes at depths $3,4,5$ fall with $N$; at depths $0,1,2$
they are within the exact floor at every octave measured, so no decay
exponent can be fitted there at all. Whether the exponents differ by
depth is therefore not decided by the range accessible here. What
Proposition~\ref{prop:scaleinv} does settle is that whatever produces
the decay, it is not the arithmetic excess that produces the size.

None of this threatens $C(N)=o(N)$: the mask is lower order under every
parameterisation considered.

\section{The negative map}\label{sec:closures}

Every entry below was pre-registered: the decision rule was fixed in
writing before the computation ran. Route adjudications were done in
fresh context against the source papers' verbatim lemma hypotheses. The
verdicts are stated here without their supporting statistics; those
live in the repository, where each carries its own null and its own
error bar. Where a design's threshold was chosen as an effect size
rather than in standard errors of its own null, the repository records
that too, and the verdict rests on the measurement rather than on the
criterion.

\subsection{Adjudication of existing machinery (5)}

\begin{longtable}{p{0.24\textwidth}p{0.10\textwidth}p{0.56\textwidth}}
\hline
Route & Verdict & Blocking coordinate\\
\hline
\endhead
MRT \cite{MRT15} / Lichtman \cite{Li20}, shift $\to$ dilate & Blocked &
 the orthogonality factorization needs a \emph{linear} pair
 constraint ($h = m-n$); the dilate constraint
 $m'u - mu' = N(m'-m)$ is bilinear. The $h$-average is a translation;
 the $k$-average is a dilation, with no diagonalizing character
 family.\\
Tao 2016 \cite{Tao} entropy decrement, $k$-averaged & Blocked $\times 3$ &
 (i) no $k$-analog of approximate affine invariance; (ii) the sampling
 prime enters the phase multiplicatively, so the sparsification fails
 and the surviving input is the target itself; (iii) the saving is
 triple-log, below spec.\\
Lichtman rerun in dilate coordinates \cite{Li23} & Blocked &
 the congruence coupling is a genuine isomorphism but powers only the
 typical-set restriction; the decoupling blocks at the same bilinear
 coordinate, and the uniformity slot would need an $EH_\mu$-grade
 input (circular).\\
Dirichlet-polynomial fourth moment / Perron & Blocked &
 every log-saving mean-value pillar assumes coefficient
 multiplicativity in its own variable; $\mu(N-u)$ has none. Unfolding
 by Perron costs $T \gtrsim N^{1-o(1)}$; opening the fourth moment
 reproduces the binary correlation.\\
Partial slices & Partial &
 the type-I slice is already consumed; the \emph{$N$-averaged}
 theorem is provable but lands in exceptional-set territory, which
 Huang--Li cannot consume; no nontrivial fixed-$N$ slice exists.\\
\hline
\end{longtable}

\noindent
\textbf{The common obstruction.} $\mu(m)\mu(N-mk)$ couples its
variables simultaneously through the product $mk$ and the difference
$N-mk$; the pair constraint is bilinear and is diagonalized by no
single character family, additive or multiplicative, while each
route's decisive lemma consumes exactly that diagonalization as a
hypothesis.

\subsection{Technique designs, kill-tested (9)}

\begin{longtable}{p{0.05\textwidth}p{0.30\textwidth}p{0.55\textwidth}}
\hline
\# & Design & Result (pre-registered rule applied)\\
\hline
\endhead
K1 & multiplicative Fej\'er kernel on the exact dilation ladder &
 \textbf{open}. On the type-II field the design is not computable: the
 $\sqrt N$ cut empties all but thirteen of the sixty-three orbit
 columns, so what was measured is a fifth of the orbit the design
 names. On the untruncated field, where the whole orbit is live, the
 statistic lands in the band the pre-registration reserved for
 ``repeat at a second $N$ before deciding''.\\
K2 & manufactured pair congruence (determinant/Kloosterman) &
 \textbf{dead}: congruent pairs are statistically $h$-blind, and the
 verdict is quantitative --- the design would need several orders of
 magnitude more than its own detection floor.\\
K3 & Wishart / operator moment method &
 \textbf{dead}: the second, third and fourth traces sit on the Wishart
 null. No sub-Wishart surplus to fund the exchange.\\
K4 & $N$-average descent &
 \textbf{dead}: the dual field is $\delta$-blind and
 $m$-uncorrelated.\\
R1 & zero-spectrum visibility (explicit formula) &
 \textbf{dead} on the measurement, which sits below its own null. The
 quoted precision of that null is not supportable --- its spread was
 estimated from six draws --- so no exclusion interval is claimed.\\
R2 & determinant / Kloosterman phase &
 \textbf{dead} on the measurement, which sits below its control. Its
 two criteria are uncalibrated and neither decides; and the
 coherent-gain arm is not blind but sits about $2.8$ standard errors
 above its own null, under a bar that was set at $5.4$.\\
R3 & character transform of the $k$-average &
 \textbf{dead by analysis}: for $K \le N^{1/3}$ the transform deposits
 every component into thin progressions (moduli $\ge N^{2/3}$);
 Parseval forbids a statistical gain.\\
R4 & divisor switch &
 \textbf{dead}: see \S\ref{sec:R4}.\\
R5 & the circle method applied directly to $C(N)$ &
 \textbf{dead, zero margin}: both bills lie at or above the trivial
 bound, and the cap on the pointwise route is Parseval
 (Proposition~\ref{prop:E}).\\
\hline
\end{longtable}

\begin{remark}[what a null verdict costs]
A kill-test that does not fire is evidence of absence only against a
stated floor, and a threshold chosen as an effect size is not a
threshold. Three of the rows above have been restated on that basis:
K1 is reopened because its design was not computable on the field it
was run on; R1's and R2's verdicts are retained but their quoted
precisions are withdrawn. In each case the DEAD direction, where it
stands, stands on the measurement --- never systematically below the
null --- and not on the count of flagged levels.
\end{remark}

\subsection{Representation classes (3, plus one open)}

\begin{longtable}{p{0.10\textwidth}p{0.25\textwidth}p{0.55\textwidth}}
\hline
Class & Test & Result\\
\hline
\endhead
C-I abelian & rational-peak energy vs mask-null &
 \textbf{closed}: the abelian spectrum is mask-exact.\\
C-II inverse domain & special-frequency excess in the
 modular-inverse re-indexing &
 \textbf{closed on verification}: fired under eight-draw nulls,
 collapsed under sixty-four-draw nulls, permutation null concurring,
 second-$N$ replication empty.\\
C-IV manufactured modularity & approximate Fricke law for
 $\Phi(z) = \sum t_m e(mz)$ &
 \textbf{closed}. A true modular form drives the defect to
 $\approx 0$, so absence is meaningful.\\
C-III Motohashi type & (no finite test) &
 \textbf{open}, needing exactly the three items of \S\ref{sec:c3}.\\
\hline
\end{longtable}

\subsection{R4 in detail: the divisor switch does not
localize}\label{sec:R4}

Applied to the dilate field over its \emph{full} ranges the switch
gives an exact identity
\[
  \sum_{k\ge1}\ \sum_{m:\,mk\le N-1}\mu(m)\mu(N-mk)
  = \sum_{u<N}\mu(N-u)\sum_{m\mid u}\mu(m)
  = \mu(N-1),
\]
verified by brute force at six values of $N$. This is perfect
cancellation: $O(1)$ for a double sum of $\sim N\log N$ terms, far
beyond square root, and it is the strongest cancellation found
anywhere in this work.

E1 imposes two restrictions that make the inner divisor sum
incomplete: the type-II cut $m > \sqrt N$ and the dyadic band
$k \sim K$. The kill-test asked whether any of the cancellation
survives, by measuring block sums $S_B(j) = \sum_{k\in\text{block}}D(k)$
against the $B$-independence that Conjecture~\ref{conj:L} predicts.

\emph{The statistic, stated.} Two normalisations of a block sum both
equal $1$ under independence, and they disagree; we use the
unweighted mean of per-block ratios,
\[
  r(B) \;=\; \Bigl\langle\,
    S_B(j)^2 \Big/ \sum_{k\in\text{block }j}\mathrm{supp}(k)
  \,\Bigr\rangle_j ,
\]
whose $B=1$ baseline is $1$ under exact square-root cancellation on the
surviving support. (The ratio-of-sums normalisation gives a $B=1$
baseline of $1.78$, because the five smallest $k$ carry a third of
$\sum_k\mathrm{supp}(k)$; it is reported in the repository and is not
used here.)

\emph{Result: none of the cancellation survives.} Over the entire band
on which the type-II field is non-empty --- $k<\sqrt N$, since
$m>\sqrt N$ forces it --- $r(B)$ reads $1.024,\,1.017,\,1.016,\,1.024$
at $B=1,2,4,8$: flat. It degrades only at block sizes where its own
sampling spread swamps it. The sharper diagnostic, the lag-1
autocorrelation of $D(k)/\sqrt{\mathrm{supp}(k)}$, reads $+0.0055$ at
$N=10^8$ against a $400$-draw permutation null of standard deviation
$0.0119$ --- $+0.5$ standard errors, dead zero, and if anything mildly
positive rather than the negative coherence a surviving residue would
require.

\emph{The mirror.} Switching the banded $L^1$ sum gives
$\sum_{k\sim K}D_{\text{full}}(k)
 = \sum_{u<N}\mu(N-u)\sum_{m\mid u,\ u/2K<m\le u/K}\mu(m)$, where
$m \asymp N/K \ge N^{2/3}$: the surviving M\"obius sits on the
\emph{long} variable, the exact opposite of the assignment that makes
Theorem~\ref{thm:A} work. The switch is not a technique that happens
to fail on the supply side; its single requirement --- M\"obius on the
short variable --- is precisely what the type-II cut forbids.

\subsection{C-III: what it needs}\label{sec:c3}

C-III is the one representation class with no finite test, and it is
open. Three requirements are identified; two are settled here against
the natural construction, and the third is external.

\paragraph{(1) A legitimate transform --- closed for changes of
variable.} Put $A = N-a$, $B = N-b$, so $A = mk$ and $B = m'k$; then
$m'A - mB = 0$ exactly. In centered coordinates the dilate family is
therefore the \emph{pencil of lines through the origin}, while the
shift family $B = A - h$ is a family of \emph{parallel} lines. A
pencil is characterised by its vertex, here a finite point; a parallel
family is a pencil whose vertex is at infinity. Affine maps send
finite points to finite points, and since $\mu$ lives on $\mathbb Z$
the available changes of variable are exactly the integral affine
ones. So no change of variable carries one family to the other. This
proves what a direct probe had only measured: the obvious re-indexing
is circular, the off-diagonal computed directly and via that
re-indexing agreeing exactly. What remains is the summation-formula
class, blocked in general by the roughness of the outer weight $\mu$,
with every named candidate in it already closed.

\paragraph{(2) A classification covering the type-II region ---
settled against the natural bookkeeping.} That bookkeeping needs the
M\"obius-side $a$ to satisfy $a \le y^{O(1)} = M^{o(1)}$. Measuring the
absolute Heath--Brown weight of the identity of level $J$ with cut
$z = M^{1/J}$, the overwhelming majority of it sits outside any region
of the form $a \le M^{\eta}$ with $\eta$ small, and the fraction
increases with $M$. We do not quote the fraction to three decimals:
its value depends on conventions the source statement leaves open ---
the rounding of $z$ alone moves the $J=8$ entry by $0.017$ --- and two
independent reconstructions of the weight differ by about one percent.

The obstruction is parameter-independent, and that is what carries the
paragraph. The identity with cut $z$ needs $z^J \ge x$, and its $j$-th
term has $a \le z^j$, which at $j = J$ is $x$ for every admissible
$(z,J)$ --- while the weight concentrates in exactly those high-$j$
terms: the top-$j$ share is about $0.82$ at $J=3$ and the
$j\in\{6,7,8\}$ share about $0.85$ at $J=8$, in both reconstructions.
A repair would have to bound
$\sum_a\sum_b \alpha(a)\beta(b)\mu(N-abk)$ with $\alpha$ rough and $a$
long, where every individual piece is trivial and all the content is
cancellation across $a$ --- a type-II estimate for $\mu(N-\cdot)$,
i.e.\ the wall. Every identity decomposing $\mu$ produces such a term;
one whose every piece had either a long free variable or a
divisor-structured rough coefficient would dispose of the parity
obstruction. \textbf{Completing the construction and breaking the wall
are the same task.}

\paragraph{(3) Quantitative averaged Chowla.} Shift-averaged binary
$\mu\mu$ correlations at a \emph{fixed} log-power saving, against a
best known of $(\log)^{1-c}$. This is a named external open problem
and is where C-III stands.

That third item sets the character of the obstruction. The
adjudication's central finding was that the \textbf{dilate} average
admits no diagonalizing character family; but the \textbf{shift}
average is precisely the home ground of the
Matom\"aki--Radziwi\l\l--Tao machinery. If a legitimate transform
converting one into the other exists, what remains is quantitative
strength inside an active area rather than the absence of any coupling
surface.

\section{What the closures constrain}

Any transform or invariance $T$ that linearizes the pair constraint at
a cost of at most a log power must satisfy all of the following.

\begin{enumerate}
\item \textbf{(K2)} $T$ cannot manufacture its congruence externally:
  collapse structure pays only when it arises inside an intrinsic
  average that is already present.
\item \textbf{(K3)} $T$ cannot bootstrap from the field's own moments:
  every moment consumes higher $\mu$-correlations and the field
  carries no sub-Wishart surplus.
\item \textbf{(K4)} $T$ cannot descend from the $N$-average: its
  linearization is load-bearing and the $k$-average holds no shadow
  of it.
\item \textbf{(R4)} $T$ cannot be the divisor switch or a relative:
  the switch's cancellation is a property of \emph{complete} divisor
  sums, and the type-II cut makes every divisor sum incomplete.
\end{enumerate}

\noindent
The constraint that K1 was to have supplied --- that $T$ cannot act
through divisibility sub-sums --- is \emph{not} asserted here, because
K1 is reopened above. Together with the round-2 findings ---
characters merely relocating the difficulty into thin progressions,
determinant phases not reaching their bar --- the field offers no
coupling surface in any direction that has an existing mathematical
name. In the automorphic world, shifted convolutions
$\sum a(n)a(n+h)$ are controlled because the coefficients $a(n)$
\emph{come with} a spectral representation; $\mu$ has none, and every
route above failed exactly where it tried to borrow one. The technique
to be created is a spectral (or equivalent structural) representation
of $\mu$-pairs themselves, not a projection onto existing spectra ---
and its construction is a purely theoretical act, beyond what
kill-testing can reach.

\section{The margin, and where the difficulty is}

The requirement constrains every $N$, so the figure to quote is the one
at the extreme, not at a typical $N$. With $\max|C|\approx
a_n\sqrt{V(N)}$ and $a_n$ the Gumbel location for the number of even
$N$ below the point,
\[
  \frac{N}{\max_{N\le X}|C(N)|}\;\approx\;
  \frac{\sqrt N}{a_n\sqrt{\AAA\log N}},
\]
which is $10^{4.5}$ at $N=10^{12}$ and $10^{22.9}$ at $N=10^{50}$.
Measured, $\max|C|/N$ falls from $0.114$ in the octave at
$3\cdot10^4$ to $0.0101$ in the octave at $1.6\cdot10^7$, so the margin
at the top of the computed range is a factor $99$, and the fitted decay
$\max|C|/N \sim N^{-0.43}$ extrapolates to a factor near $250$ at
$N=10^8$. The requirement is not remotely tight.

Placing the two standard estimates in the same units makes the shape of
the difficulty explicit. The target is $N(\log N)^{-A}$. The trivial
bound sits a factor $(\log N)^{A}$ above it; Cauchy--Schwarz gives
$N\sqrt{6(\log N-1)/\pi^2}$, a factor $(\log N)^{A+1/2}$ above it; and
the truth sits a \emph{power of $N$ below} it. \textbf{The entire
difficulty is a power of $\log N$}, and it is not the size but the
proof that is missing --- by Theorems~\ref{thm:D} and~\ref{thm:Dprime}
the one classical route to proving it is closed over its whole design
space.

\begin{remark}
Everything measured here is at $N\le1.6\cdot10^7$, with two arms at
$N\approx10^8$. The no-go theorems of Section~\ref{sec:demand} are
asymptotic and acquire content only near $N\approx10^{480}$; nothing
measured here constrains that range, and nothing there is contradicted
by these measurements.
\end{remark}

\section{What this version does not claim}\label{sec:notclaimed}

Stated so the gap is visible rather than inferred, and read against the
claim list of \S\ref{sec:claims}.

\begin{itemize}
\item \textbf{No law for $C(N)$ is conjectured here.} An earlier draft
  stated a conjecture of the form
  $C(N) = m(N) + \sqrt{V(N)}\,G(N)$ with $G$ Gaussian, supported by
  four measurements of the bulk, the tail, and the phase content in
  $\log N$. Re-verification found that the bulk and tail figures are
  not reproducible under the cell index and pooling the text specifies,
  that the phase-content measurement is reproduced by a coin and so by
  Lemma~\ref{lem:coin} is not measuring $\mu$, and that the mask's
  quoted significance was overstated. The conjecture may well be true;
  the evidence offered for it was not sound, and it is withheld until
  the statistics it rests on are defined and re-measured.
\item \textbf{No level or rate for $\rho$ is quoted.} By
  Lemma~\ref{lem:coin} the centred estimator cannot distinguish $\mu$
  from a coin, so any such figure calibrates nothing.
\item \textbf{No decay exponent for the mask is quoted.} At the three
  shallowest depths the amplitude does not clear the exact floor at any
  scale measured, so nothing can be fitted there.
\item Conjecture~\ref{conj:L}'s non-E1 stamps are single-witness and are
  cited as motivation, not as verified measurement.
\end{itemize}

Each of these is an open question rather than a retraction of a
mechanism, and each is carried forward in the repository's continuation
notes.

\section{Methodology}

Four rules were followed throughout, and they are the reason we believe
the negative results.

\begin{enumerate}
\item \textbf{Pre-registration.} Every design's decision rule was
  written before its computation ran, including the rule that would
  refute the hypothesis under test.
\item \textbf{Adversarial review in fresh context}, against the source
  papers' verbatim lemma hypotheses rather than against summaries of
  them, and repeated: every numbered statement here has been
  re-verified by a second reviewer working from the statements alone,
  without access to the first reviewer's findings or code.
\item \textbf{Nulls before thresholds.} A threshold means nothing until
  the spread of the quantity it judges has been measured, and a null
  must preserve the structure of the field it is a null for. A null
  built by shuffling a correlated field is not a null for a statistic
  that reads its correlation.
\item \textbf{Weights and fields before comparisons.} Two summaries of
  one object are not comparable until each one's weight is stated
  (Remark~\ref{rem:scale}), and a measured figure is not interpretable
  until the range it was measured on is stated. Both rules were adopted after the
  corresponding mistakes were made and caught.
\end{enumerate}

\noindent
Three traps are worth recording because they are traps and not slips.

\begin{itemize}
\item \emph{The $(q,N)>1$ main-term trap.} In the proof of
  Theorem~\ref{thm:A}, assigning a main term $T_m/\varphi(q)$ to a
  residue class with $(q,N)>1$ --- which is degenerate, its true
  contribution being $O(\log^2 N)$ --- shifts the density by exactly
  $N/\varphi(N)$. Carried into the $\log k$ branch, that error
  produces an apparent \emph{refutation of $EH_\mu$}. This is the most
  plausible false positive we met.
\item \emph{Nulls estimated from too few draws} inflate maxima across a
  family of tests; one representation-class ``hit'' survived
  eight-draw nulls and died under sixty-four-draw nulls plus
  replication, and one kill-test's quoted precision rests on a
  six-draw spread.
\item \emph{Local factors evaluated at the wrong modulus.} A
  reconstruction of the wall's excess divided by $W(N)$ where the
  definition calls for $V(N)$; the two differ by exactly $\AAA(N)$, the
  object Proposition~\ref{prop:V} is about.
\end{itemize}

\section{Reproducibility}

All measurements run on a laptop with Python and numpy. The
repository's \texttt{PROVENANCE.md} maps every numbered statement in
this paper to the code that verifies it and the result file it was
read from. One-shot verification of the core corpus is
\texttt{python code/verify/verify\_all.py}. The consistency gate, which
runs every consistency checker and exits nonzero if any fails, belongs
to the program's own record rather than to this paper.

\section{Summary}

\begin{itemize}
\item The demand side of the Huang--Li reduction is closed: one half
  is now an unconditional theorem, the other half is equivalent to the
  conclusion, and the interior of the design space between them is
  empty.
\item The supply side is closed to measurement: the object is
  featureless in every direction with a name, and the pre-registered
  closures say precisely which directions were checked and why each
  fails. One of them, K1, is reopened here.
\item The wall has an exact second moment, an exact aggregate identity,
  an exact closed form for the fluctuation of any cell mean, and a
  deterministic mask that is detected against that floor at every scale
  measured. Its error bar falls like $(\log N)^{-1/2}$ and not like a
  count.
\item The natural conditional route does not close: the coefficient
  amplifying $S(h)$ grows like $N/\log N$ and is nonnegative, so no
  bound on $|S(h)|$ gives $\mathrm{Var}\,C\sim V$; only signed
  cancellation across $h$ would, and nothing in the literature supplies
  it.
\item Net progress toward Goldbach is zero. The contribution is the
  map, the exact facts, one unconditional theorem that removes a
  Goldbach-neutral demand, one conjecture that any future construction
  must reproduce, and one defect report for the source paper.
\end{itemize}

\begin{thebibliography}{9}
\bibitem{HL} Jing-Jing Huang and Huixi Li, \emph{On the connection
between the Goldbach conjecture and the Elliott--Halberstam
conjecture}, arXiv:2005.03811v2 [math.NT], 2022.
\bibitem{ThmA} \emph{An unconditional bound for a M\"obius-weighted
correlation sum in fixed residue classes}, companion note,
\texttt{paper/theorem\_A.tex} in this repository.
\bibitem{Pan} Cheng-Dong Pan, \emph{A new attempt on Goldbach
conjecture}, Chinese Ann. Math. \textbf{3} (1982), 555--560.
\bibitem{Tao} Terence Tao, \emph{The logarithmically averaged Chowla
and Elliott conjectures for two-point correlations}, Forum Math. Pi
\textbf{4} (2016).
\bibitem{Li20} Jared Duker Lichtman, \emph{Averages of the M\"obius
function on shifted primes}, arXiv:2009.08969.
\bibitem{Mir49} L. Mirsky, \emph{The number of representations
of an integer as the sum of a prime and a $k$-th power free
integer}, Amer. Math. Monthly \textbf{56} (1949), 17--19.
\bibitem{MRT15} Kaisa Matom\"aki, Maksym Radziwi\l{}\l{}
and Terence Tao, \emph{An averaged form of Chowla's conjecture},
Algebra Number Theory \textbf{9} (2015), 2167--2196.
\bibitem{Vau88} R. C. Vaughan, \emph{The $L^1$ mean of exponential
sums over primes}, Bull. London Math. Soc. \textbf{20} (1988),
121--123.
\bibitem{Li23} Jared Duker Lichtman, \emph{Primes in arithmetic
progressions to large moduli, and shifted primes without large prime
factors}, arXiv:2309.08522.
\bibitem{MV17} M. Ram Murty and Akshaa Vatwani, \emph{Twin primes and
the parity problem}, J. Number Theory \textbf{180} (2017), 643--659.
\end{thebibliography}

\end{document}
